% Expanded replacement for Part V only.
% Replace the former Results section in its entirety with this section.

\section{Results and Demonstrated Capabilities}

The implemented workflow was demonstrated using the Lunar Spaceport-1 ISRU case together with simple non-ISRU scenarios. The objective of the results is not to claim that the current prototype exhaustively represents lunar surface operations. Rather, the results demonstrate that the previously disconnected decisions of scenario definition, CAD selection, simulation, and visualization now produce traceable scenario-level artifacts. A named scenario can be configured in the web interface, executed through the DES workflow, saved with its inputs and outputs, compared with another completed scenario, and selected independently for Omniverse playback.

\subsection{A. Scenario Definition, Persistence, and Comparison}

The revised web interface presents mission definition as a sequence of connected decisions rather than as a set of isolated tools. The user begins from a selected or newly created scenario, activates the systems available in the imported SysML model, selects a candidate site, places the required module instances, and generates the associated routes. The interface records the multiplicity of placed systems and the route assignments of rover fleets. It then exposes the SysML interfaces, resource and power links, and equations that define the operational scenario. The DES configuration view completes this definition by exposing the numerical assumptions that govern transport, production, storage, power generation, and power consumption.

This structure makes the scenario itself a persistent engineering object. A saved scenario includes its name, selected systems, placed instances, Urban Planning outputs, route assignments, user equations, DES configuration, results, and visualization package. The user can reload a saved scenario without rebuilding its architecture or manually reentering numerical assumptions. This directly addresses a limitation of the earlier workflow, in which an interface state could lead to a result but was not consistently preserved as a reusable scenario definition.

The scenario comparison capability provides a first mechanism for design-space exploration. Completed scenarios can be selected and compared through their available time histories and final measures of effectiveness. This allows a user to assess the consequences of changes such as rover fleet size, assigned routes, processing assumptions, or power-generation settings without losing the baseline configuration. The comparison view is functional for the available result series; its aggregation and presentation will need to become more resource-agnostic as the generic scenario vocabulary expands.

\begin{figure}[htbp]
    \centering
    \includegraphics[width=0.92\linewidth]{figures/saved_scenario_comparison}
    \caption{Saved scenario comparison interface. Completed configurations can be reloaded and compared through selected result histories and final measures of effectiveness.}
    \label{fig:saved-scenario-comparison}
\end{figure}

\subsection{B. Configurable ISRU Simulation Results}

The ISRU scenario demonstrates the transition from a fixed proof-of-concept controller to a transparent operational configuration. The interface exposes the architecture and operating assumptions that determine the simulation response: the number of rovers and plant instances, rover payloads, resource routes, processing rate, storage capacities, pickup thresholds, simulation duration, solar generation, battery conditions, and module demand models. The user can therefore change the cause of an operational result rather than only observing its consequence. For example, an increase in rover capacity or fleet size changes the cadence at which regolith becomes available to the plants; a route change changes the haul distance and transport time; a processing or storage setting changes production and dispatch behavior; and a power setting changes whether the mission remains power-feasible.

The DES outputs retain the time-dependent quantities needed to evaluate these consequences. For the ISRU case, the recorded results include resource delivery and production events, plant operation and storage state, rover energy behavior, power demand, accumulated energy consumption, battery-related state, and delivery events. These outputs provide both final measures of effectiveness and histories that can be viewed through the scenario comparison interface. The use of scenario-specific configuration and result files ensures that the values shown in a comparison correspond to the exact assumptions used to generate them.

The revised ISRU engine also demonstrates architecture-level flexibility. Multiple regolith and LOX rover instances can be represented, multiple ISRU plants can operate in parallel, and each plant maintains a separate regolith input storage. Available regolith rovers are dispatched toward the least-filled eligible plant, while LOX pickup follows the configured resource route and threshold. Rover cycles include outbound and return travel, with haul distances taken from the Urban Planning route definitions. The resulting implementation remains an ISRU behavior model rather than a complete logistics physics simulation, but its mission-specific assumptions are no longer hidden from the user.

\begin{figure}[htbp]
    \centering
    \includegraphics[width=0.90\linewidth]{figures/isru_result_histories}
    \caption{Example ISRU result histories and final measures of effectiveness for a saved scenario. The displayed series are generated from the configuration defined in the interface and can be compared with another completed scenario.}
    \label{fig:isru-result-histories}
\end{figure}

\subsection{C. Initial Generalization Beyond ISRU}

The generic behavior framework was exercised with simple cargo-relay, ice-processing, and water-transfer scenarios. These cases are deliberately simpler than a full lunar-base campaign, but they demonstrate the intended generalization mechanism. Each scenario begins with a SysML architecture and its interfaces rather than with an ISRU-specific controller. The imported graph provides the modules and resource connections, while the user can inspect or override the inferred operational classes, define the routes required by transporters, and provide the equations associated with the selected roles.

The cargo-relay case demonstrates a source, transporter, storage, and consumer sequence. The ice-processing case adds a processor supplied by an ice rover and power generator. The water-transfer case demonstrates a resource transformation and downstream water storage. These examples show that the same interface can represent resource names other than regolith and LOX and that the same compilation mechanism can create source, processing, storage, transport, and power behavior from the scenario description. The result is an initial mission-description language that can express several types of resource chain without adding a new dedicated Python controller for each example.

The scope of this result should remain clear. The generic engine provides a structured starting point for new scenario types; it does not validate every physical relationship that may be required in a detailed lunar mission. Its current behavior vocabulary, equation contracts, and validation rules are intentionally constrained. Further development is required for complex resource networks, detailed queueing and loading equipment, stochastic operations, advanced unit checking, and high-fidelity power-system behavior.

\begin{figure}[htbp]
    \centering
    \includegraphics[width=0.96\linewidth]{figures/generic_scenario_examples}
    \caption{Examples of non-ISRU scenario definitions evaluated through the generic framework: cargo relay, ice processing, and water transfer. Module roles and resource flows are derived from the imported architecture and completed by user configuration.}
    \label{fig:generic-scenario-examples}
\end{figure}

\subsection{D. CAD and Omniverse Playback Capabilities}

The multi-format CAD workflow was demonstrated with native USD-family assets and with STEP/STP, STL, and OBJ submissions. For each supported input, the pipeline produces a standardized USD asset and a browser preview. It preserves source-CAD metadata separately from visualization metadata, then writes scaled physical metadata into the corresponding SysML part. This result is significant because a CAD submission is no longer only a graphical attachment: its provenance, resolved geometry, orientation, and available physical properties become traceable elements of the mission definition. The visual scale remains governed by the SysML architectural dimensions, ensuring that a source model does not silently redefine the dimensions used by Urban Planning or the scene builder.

The scenario-specific manifest closes the loop from these upstream products to Omniverse. A completed scenario provides independent scene, waypoint, manifest, terrain, and DES result files. The extension can select the corresponding package rather than loading a single global playback state. This prevents a new run from silently replacing the visual inputs of a previous scenario and enables the user to inspect several completed alternatives in the same local Omniverse environment.

The Omniverse playback demonstrates both spatial and temporal coherence. Module footprints and rover geometry are grounded on the prepared terrain mesh, routes are projected to the terrain, and slope-colored segments provide a visual diagnostic of local route conditions. During playback, the extension uses the recorded DES log to synchronize rover movement and telemetry. The interface provides overview, follow, and rover-chase cameras, allowing a user to inspect the base at different scales and follow operations along a route. The telemetry panel is most complete for the established ISRU data schema, where it reports quantities such as plant storage, production, power state, and rover information. It is synchronized with recorded results rather than streamed from a running DES, which makes it a reproducible playback capability.

\begin{figure}[htbp]
    \centering
    \includegraphics[width=0.95\linewidth]{figures/omniverse_telemetry_playback}
    \caption{Time-synchronized DES telemetry during Omniverse playback. The selected scenario is visualized on the terrain mesh while the dashboard displays recorded operational state associated with the current simulation time.}
    \label{fig:omniverse-telemetry-playback}
\end{figure}

Taken together, these results demonstrate that ECLIPSE now maintains a more continuous digital thread from scenario definition to quantitative evaluation and immersive inspection. The contribution is not that every mission can already be modeled at arbitrary fidelity. Rather, the prototype now preserves the scenario identity, operational assumptions, geometry, and visualization inputs needed to make subsequent mission iterations more traceable, configurable, and comparable.
